\documentclass[12pt]{article}
\usepackage[utf8]{inputenc}
\usepackage[T1]{fontenc}
\usepackage{hyperref}
\title{Tanahmerah Language Resources for K-12 Curriculum Development}
\author{Compiled by Claude Sonnet 5 (claude-sonnet-5) for Andr\'as Kornai}
\date{}
\begin{document}
\maketitle
\section{Introduction}
Tanahmerah (ISO 639-3: tcm; Glottocode tana1288), also called Sumuri or Sumeri, is spoken in Sumuri District, Teluk Bintuni Regency, West Papua, Indonesia. The villages named include Tofoi, Materabu Jaya, Forada, Agoda, Saengga, Tanah Merah Baru, Onar Lama and Onar Baru. The speaker count is about 500 by the Ethnologue-style figure, and one source says about a thousand. Classification is disputed. Ross (2005) and Usher (2020) treat it as an independent branch of Trans-New Guinea, and Palmer (2018) treats it as an isolate. Do not confuse it with Tabla (ISO tnm), an unrelated Sentani-branch language spoken near Jayapura that is also called Tanah Merah.
\section{Primary Educational Resources}
I found no primers, textbooks or school programmes.
\section{Linguistic and Reference Materials}
\begin{itemize}
\item Wikipedia, Sumuri language (Tanahmerah, tcm). \url{https://en.wikipedia.org/wiki/Sumuri_language}
\item Wikipedia, Tanah Merah language, which concerns the unrelated Tabla language and is listed only so it is not mistaken for tcm. \url{https://en.wikipedia.org/wiki/Tanah_Merah_language}
\item Verbix Endangered Languages page for Tanahmerah. \url{https://docs.verbix.com/EndangeredLanguages/Tanahmerah}
\end{itemize}
\section{Religious and Community Resources}
\begin{itemize}
\item Joshua Project profile of Tanahmerah/Sumeri in Indonesia. It reports no known Bible translation. \url{https://joshuaproject.net/people_groups/15243/ID}
\end{itemize}
\section{Digital and Audio Resources}
I found no audio, video or corpus resources.
\section{Curriculum Development Context}
Very little is available online. The search found no grammar, dictionary or teaching materials. A curriculum would need new documentation work with the community and local education authorities. Indonesian is the likely language of schooling.
\end{document}